%%=============================================================================
%% Inleiding
%%=============================================================================

\chapter{\IfLanguageName{dutch}{Inleiding}{Introduction}}%
\label{ch:inleiding}

Binnen moderne softwareontwikkeling vormen end-to-end (E2E) testen een onmisbaar pijler voor kwaliteitsborging. Deze testen zijn de belangrijkste methode om te garanderen dat complexe systemen blijven functioneren zoals de eindgebruiker dat verwacht. E2E-testen verifiëren namelijk de volledige gebruikersstroom van begin tot einde, waardoor het inzichten biedt in de productiekwaliteit en de schaal van problemen.
\bigskip
\\
Traditioneel vereisen E2E testen opzetten en onderhouden een continue inbreng van ontwikkelaars, hierdoor is hun kostbare tijd weggenomen van nieuwe functionaliteiten te ontwikkelen. Dus is het belangrijk om een systeem te ontwerpen dat op een zelfvoorzienende manier de productkwaliteit meet bij elke oplevering naar de GitHub repository met minimale belasting voor het ontwikkelingsteam. 
\bigskip
\\
De kwaliteit van het product dat talentguide oplevert is nu laag, zelfs in de productieomgeving voor de eindgebruiker. Om de impact hiervan te begrijpen, moet eerst duidelijk zijn wat het platform precies doet. Talentguide is een SaaS-platform dat organisaties ondersteunt bij het structureren en interpreteren van personeelsgegevens. Het platform verwerkt data zoals cv’s, functiebeschrijvingen en interne documenten en zet deze om naar een gestructureerd competentieprofiel. Daarbij worden eerst kleinere bouwstenen geëxtraheerd, zoals taken, tools, kennis en gedragskenmerken. Binnen talentguide worden deze aangeduid als TTK’s. Deze TTK’s krijgen een expertiselevel en worden vervolgens gegroepeerd onder bredere vaardigheden. Op die manier ontstaat er een profiel dat niet enkel weergeeft welke vaardigheden een werknemer bezit, maar ook op welk niveau deze vaardigheden aanwezig zijn.
\bigskip
\\
Deze profielen vormen de basis voor verdere inzichten binnen het platform. Talentguide gebruikt de gestructureerde data om werknemers te matchen met rollen, competentie gaten te identificeren en ontwikkelingsnoden zichtbaar te maken. Het platform kan ook gelijkaardige of ontbrekende competenties voorstellen om een profiel verder te vervolledigen. Omdat deze inzichten afhankelijk zijn van correcte en betrouwbare data, is de kwaliteit van de onderliggende software belangrijk. Wanneer releases fouten bevatten of gebruikersstromen niet correct werken, kan dit rechtstreeks invloed hebben op de betrouwbaarheid van competentieprofielen, validatieflows en de gebruikerservaring. Daarom is het relevant om te onderzoeken hoe een geautomatiseerde end-to-end testingworkflow kan bijdragen aan een stabieler en betrouwbaarder SaaS-platform.

%De inleiding moet de lezer net genoeg informatie verschaffen om het onderwerp te begrijpen en in te zien waarom de onderzoeksvraag de moeite waard is om te onderzoeken. In de inleiding ga je literatuurverwijzingen beperken, zodat de tekst vlot leesbaar blijft. Je kan de inleiding verder onderverdelen in secties als dit de tekst verduidelijkt. Zaken die aan bod kunnen komen in de inleiding~\autocite{Pollefliet2011}:


\section{\IfLanguageName{dutch}{Probleemstelling}{Problem Statement}}%
\label{sec:probleemstelling}

%Uit je probleemstelling moet duidelijk zijn dat je onderzoek een meerwaarde heeft voor een concrete doelgroep. De doelgroep moet goed gedefinieerd en afgelijnd zijn. Doelgroepen als ``bedrijven,'' ``KMO's'', systeembeheerders, enz.~zijn nog te vaag. Als je een lijstje kan maken van de personen/organisaties die een meerwaarde zullen vinden in deze bachelorproef (dit is eigenlijk je steekproefkader), dan is dat een indicatie dat de doelgroep goed gedefinieerd is. Dit kan een enkel bedrijf zijn of zelfs één persoon (je co-promotor/opdrachtgever).

Het grote probleem is dat er nu geen metingen zijn uitgevoerd op de productkwaliteit, waardoor er geen maatstaf is om te zien hoe groot of klein problemen zijn. Dit komt door het gebrek aan E2E testen in de huidige React/Node.js codebase. Dit leidt ertoe dat releases naar de productieomgeving breken en een slechte gebruikerservaring opleveren.

\section{\IfLanguageName{dutch}{Onderzoeksvraag}{Research question}}%
\label{sec:onderzoeksvraag}

%Wees zo concreet mogelijk bij het formuleren van je onderzoeksvraag. Een onderzoeksvraag is trouwens iets waar nog niemand op dit moment een antwoord heeft (voor zover je % %kan nagaan). Het opzoeken van bestaande informatie (bv. ``welke tools bestaan er voor deze toepassing?'') is dus geen onderzoeksvraag. Je kan de onderzoeksvraag verder %specifiëren in deelvragen. Bv.~als je onderzoek gaat over performantiemetingen, dan 

De hoofdonderzoeksvraag luidt: hoe kan een geautomatiseerd end-to-end testsysteem worden opgezet dat het product continu meet, zodat kwaliteit is verankerd in het ontwikkelingsproces en ervoor zorgt dat ontwikkelaars geen kostbare tijd moeten nemen om volledig nieuwe testen te schrijven?

\section{\IfLanguageName{dutch}{Onderzoeksdoelstelling}{Research objective}}%
\label{sec:onderzoeksdoelstelling}

%Wat is het beoogde resultaat van je bachelorproef? Wat zijn de criteria voor succes? Beschrijf die zo concreet mogelijk. Gaat het bv.\ om een proof-of-concept, een prototype, een verslag met aanbevelingen, een vergelijkende studie, enz.

Het doel van deze bachelorproef is om de kwaliteit van het eindproduct bij talentguide te verhogen door een functioneel prototype af te leveren van een geautomatiseerde testingworkflow dat op basis van minimale ontwikkelaars input robuuste testen genereert en uitvoert.


\section{\IfLanguageName{dutch}{Opzet van deze bachelorproef}{Structure of this bachelor thesis}}%
\label{sec:opzet-bachelorproef}

% Het is gebruikelijk aan het einde van de inleiding een overzicht te
% geven van de opbouw van de rest van de tekst. Deze sectie bevat al een aanzet
% die je kan aanvullen/aanpassen in functie van je eigen tekst.

De rest van deze bachelorproef is als volgt opgebouwd:

In Hoofdstuk~\ref{ch:stand-van-zaken} wordt een overzicht gegeven van de stand van zaken binnen het onderzoeksdomein, op basis van een literatuurstudie.

In Hoofdstuk~\ref{ch:methodologie} wordt de methodologie toegelicht en worden de gebruikte onderzoekstechnieken besproken om een antwoord te kunnen formuleren op de onderzoeksvragen.

In Hoofdstuk~\ref{ch:fase1} wordt een nulmeting uitgevoerd om concrete problemen te identificeren.

In Hoofdstuk~\ref{ch:fase2} wordt de requirement analyse uitgevoerd.

In Hoofdstuk~\ref{ch:fase3} wordt het prototype ontwikkeld.

In Hoofdstuk~\ref{ch:fase4} wordt het prototype geëvalueerd.

In Hoofdstuk~\ref{ch:fase5} wordt meetbaar de evaluatie gedaan.

% TODO: Vul hier aan voor je eigen hoofstukken, één of twee zinnen per hoofdstuk

In Hoofdstuk~\ref{ch:conclusie}, tenslotte, wordt de conclusie gegeven en een antwoord geformuleerd op de onderzoeksvragen. Daarbij wordt ook een aanzet gegeven voor toekomstig onderzoek binnen dit domein.